\section{Make the figure sequence support the argument}
\label{sec:figures}

Readers often meet the evidence through figures before they read the full
paper. The figure sequence should make the same argument in shorter form. A
reader who scans the title, abstract, and figures should understand the problem
and the main result. If the paper proposes an explanation, the figures should
show the evidence for it.

\subsection{Use Figure 1 to summarize the paper}

Figure~1 should identify the new intervention and show its main effect. If the
paper makes a mechanism claim, the figure should also show the evidence that
motivates it. Most first figures pair a method diagram with one headline
result. Add a task image or mechanism sketch only when it supplies necessary
context.

The example papers make different choices because their claims differ.
\emph{Action Chunking} places its two interventions in one conceptual frame.
\emph{DPPO} joins the nested decision process to robot tasks and the properties
it studies. \emph{Much Ado About Noising} uses a component taxonomy to reject
prevailing explanations and introduce a minimal replacement. \emph{OGPO}
places the algorithm beside sample-efficiency curves, then uses its final
panel to motivate the mechanism study. In each case, the first figure reflects
the paper's main claim rather than a standard figure template.

Design Figure~1 after the claim--evidence map is stable. If the figure requires
a claim that is not tested in the paper, change the figure or collect the missing
evidence. If the figure cannot state the contribution without several
paragraphs of caption, the contribution hierarchy may still be unsettled.

\subsection{Use later figures to answer distinct questions}

A later figure should resolve one question left open by Figure~1. Useful figure
types include:

\begin{description}[leftmargin=*,style=nextline]
    \item[\textbf{Concept figure.}]
    Makes a theoretical object, information flow, or mechanism visible.
    \item[\textbf{Controlled comparison.}]
    Changes one disputed component while holding the relevant alternatives
    fixed.
    \item[\textbf{Headline benchmark.}]
    Establishes the main comparison across representative regimes.
    \item[\textbf{Mechanism figure.}]
    Shows trajectories, distributions, gradients, or intermediate
    computations that bear on an explanation.
    \item[\textbf{Qualitative sequence.}]
    Displays an ordered success or failure with the event of interest marked.
    \item[\textbf{Ablation.}]
    Removes or varies one design decision and reports the resulting change.
\end{description}

Do not ask one figure to answer every question. A benchmark grid becomes hard
to read when it also contains the method diagram, ablation, and mechanism
claim. Separate figures let the prose distinguish performance from
explanation.

\subsection{Keep colors and labels consistent}

Assign colors by meaning and keep those assignments stable. Use neutral gray
for context and most baselines, then reserve one or two accent colors for the
method and the main alternative. Learning curves should state the uncertainty
shown and the number of seeds or trials. Direct labels are often easier to read
than a distant legend. Panel headings should name the comparison, not merely
``left'' and ``right.''

Diagrams need a visible direction of travel. Give each arrow one meaning and
state that meaning in the caption. Task images can ground an abstract diagram,
but they should show the actual task or failure under discussion. A decorative
image does not supply evidence.

Inspect every figure at its final column width. Check every label and
annotation at that size. Lines must remain distinguishable on a printed page,
including in grayscale and under common forms of color-vision deficiency.

\subsection{State the evidence in each caption}

Start the caption with the figure's claim or purpose. Walk through the panels
in reading order. Then report the comparison and main observation before
giving the narrowest interpretation the figure supports. Point to setup
details when they are not visible nearby.

A useful caption order is:

\begin{enumerate}[leftmargin=*]
    \item claim or purpose;
    \item panel-by-panel setup;
    \item comparison and metric;
    \item main observation;
    \item narrow interpretation or caveat;
    \item pointer to complete details.
\end{enumerate}

The caption and results paragraph should agree. A caption cannot say that an
ablation ``demonstrates the mechanism'' while the text calls it suggestive.
Define nonstandard symbols and uncertainty either in the caption or in the
adjacent prose. The reader should not need to search the appendix for the
meaning of an error band.
